Un ull hipermetrop no és capaç d'enfocar objectes propers a la retina, sinó que els enfoca darrere d'aquesta, i per això els veu borrosos. L'Anna no és capaç d'enfocar bé els objectes que són més a prop de 70~cm i necessita unes ulleres amb lents convergents per a llegir.

\begin{apartat}{1,25}
L'Anna normalment situa el llibre a 35~cm i la imatge creada per la lent ha de ser a 70~cm perquè la pugui enfocar correctament. Calculeu la distància focal de la lent correctora i la seva potència.
\end{apartat}

\begin{apartat}{1,25}
Calculeu on es formarà la imatge d'un objecte de 10~cm d'alçària situat a 25~cm davant d'una lent convergent de 20~cm de distància focal. Calculeu, també, la mida de la imatge i esmenteu les seves característiques (més gran o més petita, real o virtual, dreta o invertida). Dibuixeu a la quadrícula de sota el diagrama de raigs amb la lent, l'objecte i la imatge.
\end{apartat}

\figura[0.6]{fig1.png}
